\documentclass[a4paper,12pt]{article}
\usepackage[utf8]{inputenc}
\usepackage{amsmath}
\usepackage{parskip}
\usepackage{geometry}
\geometry{margin=1in}

\title{Propuesta de Tesina}
\author{}

\begin{document}

\maketitle

\begin{flushleft}
\textbf{Postulante:} Nombre del Estudiante \\
\textbf{Director:} Nombre del Director \\
\textbf{Codirector:} Nombre del Codirector \\
\end{flushleft}

\section{Situación del postulante}
Indique las materias que le falta aprobar para finalizar la carrera. Indique si está trabajando o realizando otra actividad que podría afectar su dedicación a la realización de la tesina, y cuántas horas por semana le dedica a esta actividad.

\section{Título}
Título de la propuesta

\section{Motivación y Objetivo General (máx. 1 pág.)}
Explique el problema o situación de referencia en el que se desarrolla la propuesta o los interrogantes en el campo disciplinario a los que la propuesta se dirige. Desarrolle la importancia e impacto de los objetivos y el conocimiento que se generará. En esta sección no es necesario describir las tareas específicas que se realizarán (para eso, ver Objetivos específicos).

\section{Fundamentos y estado de conocimiento sobre el tema (máx. 2 pág.)}
Escriba una breve introducción general al tema y cite y comente las mayores contribuciones en el tema específico, incluyendo bibliografía actualizada.

\section{Objetivos específicos (máx. 1 pág.)}
Enuncie de manera clara las metas concretas a alcanzar en el marco de la tesina.

\section{Metodología y Plan de Trabajo (máx. 2 pág.)}
Se recomienda estructurar esta sección en función de los objetivos específicos.

Planteo de la hipótesis a analizar en cada objetivo o sección del proyecto.

Actividades propuestas y metodología a usar en cada una de ellas.

Resultados que se esperan obtener o metas a cumplir y cómo se evaluarán los resultados.

Trate de evaluar los potenciales problemas y limitaciones de la metodología y técnicas propuestas y en lo posible proponer alternativas.

\subsection{Cronograma de Trabajo}
Se presentará un cronograma de trabajo con las tareas desagregadas y los tiempos estimados que consumirán.

\section{Referencias}
Referencias bibliográficas.

\end{document}
 
